%
% 28-erfahrung.tex -- Erkenntnisse
%
% (c) 2026 Prof Dr Andreas Müller
%
\subsection{Erkenntnisse}
Aus den langwierigen Rechnungen für das Integrationsproblem in
$\mathbb{Q}(x,\vartheta)$ mit $\vartheta=\!\sqrt{ax^2+bx+c}$
können wir einige Beobachtungen ableiten,
die auch für das Integrationsproblem für beliebige elementare
Funktionen gelten könnten.
\begin{enumerate}
\item 
Das Integrationsproblem spielt sich in einem Körper ab, der als
Erweiterung der rationalen Funktionen $K=\mathbb{Q}(x)$ um ein weiteres 
Element $\vartheta$ erweitert worden ist.
Der Integrand ist also eine rationale Funktion $f\in K(\vartheta)$, sie
kann als rationale Funktion 
\[
f
=
\frac{p(\vartheta)}{q(\vartheta)}
\qquad
\text{mit}
\qquad
p(\vartheta),q(\vartheta)
\in
\mathbb{Q}(x)[\vartheta]
=
K[\vartheta]
\]
geschrieben werden.
\item
Mit der Polynomdivision kann man das Integral in einen polynomiellen
und einen rationalen Teil aufteilen:
\[
\int f
=
\underbrace{
\int \text{Polynom in $\vartheta$}
}_{
\displaystyle\text{polynomieller Teil}
}
+
\underbrace{
\int\frac{\text{Rest}(\vartheta)}{q(\vartheta)}
}_{
\displaystyle\text{rationaler Teil}
}.
\]
Die Koeffizienten der Polynome sind rationale Funktionen in $K$.
\item
Die Partialbruchzerlegung kann dazu verwendet werden, den rationalen
Teil auf eine kleine Menge einfacher Grundintegrale für diese
Art von Erweiterung zu reduzieren.
Im Kapitel~\ref{chapter:ratint} war dies die Methode von Hermite.
\item
Es braucht ein spezielle Methode, um die Grundintegrale zu bestimmen.
Im Kapitel~\ref{chapter:ratint} war dies die Methode von Rothstein-Trager.
\item
Die entstehenden Stammfunktionen haben eine gemeinsame Struktur.
Sie bestehen aus einer rationalen Funktion von $\vartheta$ und
einer Linearkombination von Logarithmusfunktionen, deren Argumente
lineare Funktionen von $\vartheta$ mit Koeffizienten in $K$ sind.
\end{enumerate}
Im nächsten Abschnitt soll gezeigt, werden, dass die in Beobachtung~5
festgestellte allgemein Form der Stammfunktion tatsächlich 
immer dann vorliegt, wenn die Stammfunktion elementar ist.
Sie ist unter dem Namen Liouville-Prinzip bekannt.



